\section{Compact Systems and Extensions} \label{sec:compact}

This section proves part (K) of Lemma \ref{l:Extensions}, in Theorem \ref{t:compactExtension}: a compact extension of a CSZ factor is CSZ.  
Taking the factor to be trivial, it shows that every compact system is CSZ.  
In \S\ref{sub:rootsOfUnity} we add Theorem \ref{t:compactSystem}, which is \emph{not} used in the proof of the main theorem.  
It treats the compact systems whose eigenvalues are all roots of unity, by a direct argument, and for them it gives more: 
a lower bound of $\int f ^{t+1}\;dm - \epsilon$, attained on a syndetic set of indices.  
Throughout, let $\mathcal C = \mathcal C_{b,D} = \{0=k_0<k_1<\cdots\}$ be a classical Cantor set, with $0\in D$, and $\delta = \lvert D\rvert$.  


\subsection{Compact extensions} 

 
\begin{theorem} \label{t:compactExtension}
Let $\mathcal D$ be a factor of $(X, \mathcal A, m, T)$ which is CSZ, and let $\mathcal E$ be a compact extension of $\mathcal D$.  
Then $\mathcal E$ is CSZ.  That is, for all integers $t \geq 1$, all $i \geq 0$, and all 
non-zero $0\leq f\leq 1$ in $L ^{\infty}(\mathcal E)$, we have 
\begin{equation} \label{e:compactExtensionConclusion}
\liminf_{N-M\to\infty}  \Expectation_{M\leq n<N} 
\int \prod _{s=0}^{t} T ^{s b ^{i} k_n} f\;dm >0. 
\end{equation}
\end{theorem}

The proof follows Furstenberg, Katznelson and Ornstein \cite{MR670131}, in the form presented by Moreira \cite{MoreiraNotes}*{\S 8.5}.  
We single out two features of that argument which make it the right one for our purposes.  
First, it works directly with the uniform averages that we require, whereas the van der Waerden argument of 
Furstenberg \cite{MR0498471} is naturally phrased with ordinary Ces\`aro averages.  
Second, and decisively, the return time that it produces is a \emph{translate} $n+a$ of the averaging parameter $n$, 
by an element $a$ of a fixed finite set, whereas the van der Waerden argument produces a \emph{dilate} $sn$, with $s$ 
ranging over a finite set.  The integers are closed under dilation, so for Szemer\'edi's theorem either argument serves.  
A Cantor set is not closed under dilation, but it is, in the following sense, closed under the translations we need: 
if $a \in \mathcal C\cap [0,b ^{i})$ and $n \in b ^{i}\mathcal C$, then by Lemma \ref{l:Diff}(ii), 
\begin{equation} \label{e:absorb}
n = b ^{i}k_u \quad\Longrightarrow\quad n+a = b ^{i}k_u + a = k _{\delta ^{i} u + \iota (a)} \in \mathcal C , 
\end{equation}
where $\iota (a) < \delta ^{i}$ denotes the index of $a$ in $\mathcal C$.  
That is, the translate stays in the Cantor set, and its index stays inside the block $[\delta ^{i}u , \delta ^{i}(u+1))$.  
This is why the argument below restricts the finite set of translations to the Cantor set, and the averaging parameter to a dilate of it.  

\begin{remark}
In particular, the proof below does not use the Cantor van der Waerden Theorem of \S\ref{sec:vdW}, 
which is retained as a result of independent interest.  
\end{remark}

We need this known Lemma. 

\begin{lemma} \label{l:localize} Let $\mathcal E\mid \mathcal D$ be a compact extension. 
Then for all $f\in L^2(\mathcal E\mid \mathcal D)$ and $0< \epsilon <1$, there is a set $A \in \mathcal D$ such that $m(A)>1- \epsilon $ and $f \mathbf{1}_{A}$ is conditionally compact. 
\end{lemma}

See \cite{MoreiraNotes}*{Lemma 8.4.3} for a proof.

\begin{proof}[Proof of Theorem \ref{t:compactExtension}]
We give the argument for the dilation $i=0$.  The case of a general $i$ is the same argument with $T$ replaced by $T ^{b ^{i}}$, 
and is set out in Remark \ref{r:dilation} at the end of the section.  
The hypothesis that $\mathcal D$ is CSZ enters exactly once, at \eqref{e:c0} below, and there for the same Cantor set $\mathcal C$ 
and a dilation fixed in the course of the proof.  

Fix an integer $t \geq 1$ and a non-zero $0 \leq f \leq 1$ in $L^2(\mathcal E \mid \mathcal D)$.   
We may and do assume that $f$ is conditionally compact.  Indeed, as $f \geq 0$ is not identically zero, we may choose $0<\eta <1$ so small 
that $m(A_\ast) > 1- \eta$ implies $\int_{A_\ast} f \;dm > 0$; let $A_\ast$ be as in Lemma \ref{l:localize} for this $\eta$, and replace $f$ by $f \mathbf 1_{A_\ast}$.  
This is legitimate because $0 \leq f \mathbf 1_{A_\ast} \leq f$ and all the factors in \eqref{e:compactExtensionConclusion} are non-negative, 
so a positive liminf for $f \mathbf 1_{A_\ast}$ gives one for $f$.  
Set 
\begin{equation}
g \coloneqq \mathbb E (f \mid \mathcal D), \qquad \delta _0 \coloneqq \int_X f \;dm >0 . 
\end{equation}
Now fix, once and for all, an $\epsilon >0$ small enough that 
\begin{equation} \label{e:epschoice}
c_1 \coloneqq ( \delta _0/2) ^{t+1} - t \epsilon >0 . 
\end{equation}
Thus $\epsilon$ depends only on $t$ and $\delta _0$; everything constructed below depends on $\epsilon$.  

\smallskip 
\noindent \textbf{Good pairs.}  
By conditional compactness there are $r \in \mathbb N$ and $f_1 , \ldots, f_r$ so that for almost every $x$ and every $n \in \mathbb N_0$, 
\begin{equation} \label{e:net}
\min_{1 \leq \ell \leq r} \lVert T^{n}f - f_\ell \rVert _{\mathcal E \mid \mathcal D}(x) < \epsilon /2 . 
\end{equation}
Call a pair $\{a, a'\}$ of distinct elements of $\mathcal C$ \emph{good at $x$} if 
\begin{equation}
\lVert T^{sa}f - T ^{sa'}f \rVert _{\mathcal E \mid \mathcal D}(x) < \epsilon , \qquad 0 \leq s \leq t , 
\end{equation}
and call a finite $F \subset \mathcal C$ \emph{unfortunate at $x$} if no pair in $F$ is good at $x$.  

\begin{lemma} \label{l:unfortunate}
Set $M \coloneqq r ^{t+1}+1$.  If $F \subset \mathcal C$ has $\lvert F\rvert \geq M$, then for almost every $x$, 
$F$ is not unfortunate at $x$.  
\end{lemma}

\begin{proof}
Let $Y$ be the set of full measure on which \eqref{e:net} holds for all $n$, and fix $x \in Y$.  
To each $a \in F$ associate the tuple $( \ell _0 , \ldots, \ell _t) \in [r] ^{t+1}$, where $\ell _s$ is least with 
$\lVert T^{sa}f - f _{\ell _s} \rVert _{\mathcal E\mid \mathcal D}(x)< \epsilon /2$.  
There are $r ^{t+1}$ such tuples, so if $\lvert F \rvert \geq M$ two elements $a \neq a'$ of $F$ share a tuple.  
The conditional triangle inequality then gives $\lVert T^{sa}f - T^{sa'}f\rVert _{\mathcal E\mid \mathcal D}(x)< \epsilon$ for $0 \leq s \leq t$. 
\end{proof}

\smallskip 
\noindent \textbf{A maximal unfortunate set.}  
For finite $F \subset \mathcal C$ put 
\begin{equation}
S(F) \coloneqq \{ x \colon F \textup{ is unfortunate at } x \} 
= \bigcap_{a \neq a' \in F} \ \bigcup_{0\leq s \leq t} 
\bigl\{ x \colon \lVert T^{sa}f-T^{sa'}f \rVert _{\mathcal E\mid \mathcal D}(x) \geq \epsilon \bigr\} .
\end{equation}
Conditional norms are $\mathcal D$ measurable, so $S(F) \in \mathcal D$.  Put 
\begin{equation}
Q(F) \coloneqq \bigl\{ x \in S(F) \colon x \notin S(G) \textup{ for every finite } G \subset \mathcal C \textup{ with } \lvert G\rvert > \lvert F\rvert \bigr\} . 
\end{equation}
As $\mathcal C$ is countable, it has only countably many finite subsets; so $Q(F)$ is a countable intersection of sets in $\mathcal D$, and is in $\mathcal D$.  
Singletons are unfortunate at every $x$, and by Lemma \ref{l:unfortunate} the sets $S(G)$ with $\lvert G\rvert \geq M$ are null.  
Hence almost every $x$ lies in $Q(F)$ for some finite $F \subset \mathcal C$ with $0< \lvert F\rvert < M$, namely for $F$ an unfortunate set at $x$ of maximal cardinality.  

Let $A \coloneqq \{ g > \delta _0/2\} \in \mathcal D$, and note that $\delta _0 = \int_A g + \int _{X \setminus A} g \leq m(A) + \delta _0/2$, 
so that $m(A) > \delta _0/2$.  
There are countably many finite subsets of $\mathcal C$, so we may fix one, $F$, with $m (Q(F) \cap A)>0$.  
For $x \in Q(F)$ no pair of $F$ is good at $x$, so for each pair $a \neq a'$ there is an $s$ with 
$\lVert T^{sa}f-T^{sa'}f\rVert _{\mathcal E\mid \mathcal D}(x) \geq \epsilon$; necessarily $s \in [t]$, since the case $s=0$ reads $0 \geq \epsilon$.  Let $F ^{(2)} \coloneqq \{ (a,a') \in F \times F \colon a \neq a'\}$, the ordered pairs of distinct elements.  
For each of the finitely many maps $q \colon F ^{(2)} \to [t]$, let $E_q$ be the set of $x \in Q(F) \cap A$ with 
\begin{equation} \label{e:qmap}
\lVert T^{q(a,a')a}f-T^{q(a,a')a'}f\rVert _{\mathcal E\mid \mathcal D}(x) \geq \epsilon , 
\qquad (a,a') \in F ^{(2)} . 
\end{equation}
Each $E_q$ is in $\mathcal D$, since conditional norms are $\mathcal D$-measurable, and by the previous sentence the sets $E_q$ cover $Q(F) \cap A$.  
So we may fix a map $q$ with $m(E_q)>0$, and we set $E \coloneqq E_q$.

\smallskip 
\noindent \textbf{Fixing the scale.}  
The set $F$ is a fixed finite subset of $\mathcal C$, so we may fix $i \geq 1$ with 
\begin{equation} \label{e:scale}
F \subset \mathcal C \cap [0, b ^{i}) . 
\end{equation}
This is the scale at which the argument will produce its return times, through \eqref{e:absorb}.  

\begin{lemma} \label{l:returnTime}
Let $u \geq 1$ and set $n = b ^{i}k_u$.  Let $x \in E \cap T ^{-n}E \cap \cdots \cap T ^{-tn}E$.  
Then there is an $a = a(u,x) \in F$ with 
\begin{equation} \label{e:good0}
\lVert T ^{s(n+a)}f - f \rVert _{\mathcal E\mid \mathcal D}(x) < \epsilon , \qquad 0 \leq s \leq t . 
\end{equation}
\end{lemma}

\begin{proof}
Let $a \neq a'$ in $F$ and set $s_0 = q(a,a') \in [t]$.  Since $T ^{s_0 n}x \in E$, \eqref{e:qmap} applies at the point $T ^{s_0n}x$.  
The conditional norm is equivariant, in that $\lVert T ^{v} h \rVert _{\mathcal E\mid \mathcal D}(x) = \lVert h \rVert _{\mathcal E\mid \mathcal D}(T ^{v}x)$, so that 
\begin{equation}
\lVert T ^{s_0(n+a)}f - T ^{s_0 (n+a')}f \rVert _{\mathcal E\mid \mathcal D}(x) 
= \lVert T ^{s_0a}f - T ^{s_0a'}f \rVert _{\mathcal E\mid \mathcal D}(T ^{s_0n}x) \geq \epsilon . 
\end{equation}
Thus no pair of distinct elements of $F + n$ is good at $x$.  

Now set $G \coloneqq (F+n) \cup \{0\}$.  By \eqref{e:scale} and \eqref{e:absorb}, $F + n \subset \mathcal C$, and $0 \in \mathcal C$, so $G \subset \mathcal C$.  
Moreover $n = b ^{i}k_u \geq b ^{i} > \max F$ as $u \geq 1$, so $0 \notin F+n$ and $\lvert G\rvert = \lvert F\rvert +1$.  
Since $x \in E \subset Q(F)$, the set $G$ is not unfortunate at $x$, so it contains a good pair.  By the previous paragraph that pair must involve $0$; 
say it is $\{n+a, 0\}$ with $a \in F$.  This is \eqref{e:good0}. 
\end{proof}

\smallskip 
\noindent \textbf{Conclusion.}  
Apply the CSZ property of $\mathcal D$ to the set $E$, with dilation $i$, to get  
\begin{equation} \label{e:c0}
c_0 \coloneqq \liminf_{N-M\to\infty}  \Expectation_{M\leq u<N} m (I_u)>0 , 
\qquad I_u \coloneqq E \cap T ^{-b ^{i}k_u}E \cap \cdots \cap T ^{-tb ^{i}k_u}E . 
\end{equation}
Let $B \coloneqq \{u \geq 1\colon m (I_u) > c_0/2\}$.  Splitting the average over $B$ and its complement, and using $m(I_u) \leq 1$, 
\begin{equation}
c_0 \leq \liminf_{N-M\to\infty} \frac{\lvert B \cap [M,N) \rvert}{N-M} + \frac {c_0}2 , 
\end{equation}
so that $B$ has lower uniform density at least $c_0/2$.  

Fix $u \in B$, write $n = b ^{i} k_u$, and let $x \in I_u$.  Lemma \ref{l:returnTime} provides $a = a(u,x) \in F$ satisfying \eqref{e:good0}.  
Write $d = n+a$.  Telescoping, $\prod _{s=0}^{t} T ^{sd}f - f ^{t+1}$ is a sum of $t$ terms, each of the form 
$P \cdot (T ^{sd}f - f) \cdot P'$ with $1\leq s \leq t$ and $P, P'$ products of the remaining factors, so bounded by $1$.  
Applying the conditional Cauchy--Schwarz inequality to each term, and then the conditional Jensen inequality, gives 
\begin{equation}
\mathbb E \Bigl[ \prod _{s=0}^{t} T ^{s(n+a)}f \;\Big\vert\; \mathcal D \Bigr](x) 
\geq \mathbb E \bigl( f ^{t+1} \mid \mathcal D\bigr)(x) - t \epsilon 
\geq g ^{t+1}(x) - t\epsilon 
\geq (\delta _0/2) ^{t+1} - t \epsilon = c_1 >0, 
\end{equation}
the last two steps because $I_u \subset E \subset A$, and by the choice \eqref{e:epschoice} of $\epsilon$.  
Fix an ordering of $F$, and let $a(u,x)$ be the least $a \in F$ for which \eqref{e:good0} holds at $x$.  
For each $a \in F$, the set of $x$ at which \eqref{e:good0} holds is in $\mathcal D$, as the conditional norms are $\mathcal D$-measurable; 
so $x \mapsto a(u,x)$ is $\mathcal D$-measurable on $I_u$, and it takes at most $\lvert F\rvert$ values.  
Hence there is an $a_u \in F$ for which $J_u \coloneqq \{ x \in I_u \colon a(u,x) = a_u\}$, a set in $\mathcal D$, has $m(J_u) \geq m(I_u)/\lvert F\rvert$.  
Since $f \geq 0$, 
\begin{equation} \label{e:lowerbound}
\int_X \prod _{s=0}^{t} T ^{s(n+a_u)}f \;dm 
\geq \int _{J_u} \mathbb E \Bigl[ \prod _{s=0}^{t} T ^{s(n+a_u)}f \;\Big\vert\; \mathcal D \Bigr] dm 
\geq \frac{m(I_u)\, c_1}{\lvert F\rvert} \geq \frac{c_0c_1}{2 \lvert F\rvert} \eqqcolon c_2 >0 . 
\end{equation}

Finally we convert \eqref{e:lowerbound} into a statement about the enumeration of $\mathcal C$.  
Since $0 \in D$, the first $\delta ^{i}$ elements of $\mathcal C$ are exactly $\mathcal C \cap [0,b ^{i})$, by Lemma \ref{l:Diff}(i).  
So $a_u = k _{\iota (a_u)}$ for a unique index $\iota (a_u) < \delta ^{i}$, and Lemma \ref{l:Diff}(ii) gives, as in \eqref{e:absorb}, 
\begin{equation}
n + a_u = b ^{i}k_u + k _{\iota (a_u)} = k _{v(u)} , \qquad v(u) \coloneqq \delta ^{i}u + \iota (a_u) . 
\end{equation}
Thus $v(u)$ lies in the block $[\delta ^{i}u , \delta ^{i}(u+1))$, and 
$\int_X \prod _{s=0}^{t} T ^{s k_{v(u)}}f \;dm \geq c_2$ for every $u \in B$, by \eqref{e:lowerbound}.  
Let $V \coloneqq \{ v(u) \colon u \in B\}$.  The blocks $[\delta ^{i}u, \delta ^{i}(u+1))$ are consecutive intervals of length $\delta ^{i}$ partitioning $\mathbb N_0$, 
so $u \mapsto v(u)$ is injective and $V$ meets the block of each $u \in B$ in exactly one point.  
If $\lceil P/\delta ^{i}\rceil \leq u < \lfloor Q/ \delta ^{i} \rfloor$ then the whole block of $u$ lies in $[P,Q)$, so 
\begin{equation} \label{e:count}
\bigl\lvert V \cap [P,Q) \bigr\rvert 
\geq \bigl\lvert B \cap [ \lceil P/\delta ^{i}\rceil , \lfloor Q/ \delta ^{i} \rfloor ) \bigr\rvert . 
\end{equation}
The interval on the right has length at least $(Q-P)/ \delta ^{i} -2$, which tends to infinity with $Q-P$.  
Dividing \eqref{e:count} by $Q-P$ and using that $B$ has lower uniform density at least $c_0/2$, 
\begin{equation}
\liminf_{Q-P\to\infty} \frac{\lvert V \cap [P,Q)\rvert}{Q-P} \geq \frac{1}{\delta ^{i}}\cdot \frac {c_0}2 . 
\end{equation}
The integrand $\int_X \prod _{s=0}^{t} T ^{s k_{v}}f \;dm$ is non-negative for every index $v$, and is at least $c_2$ for $v \in V$.  Therefore 
\begin{equation}
\liminf_{Q-P\to\infty} \Expectation_{P \leq v <Q} \int_X \prod _{s=0}^{t} T ^{s k_{v}}f \;dm 
\geq \frac{c_0c_2}{2 \delta ^{i}} >0 . 
\end{equation}
This is \eqref{e:compactExtensionConclusion} with $i=0$, and so completes the proof. 
\end{proof}

\begin{remark} \label{r:dilation}
To obtain \eqref{e:compactExtensionConclusion} for a general dilation $i_0 \geq 0$, run the proof above with $T$ replaced by 
$\widetilde T \coloneqq T ^{b ^{i_0}}$ throughout.  This is legitimate: $\mathcal E$ is a compact extension of $\mathcal D$ for $\widetilde T$, 
by Lemma \ref{l:powers}(1).  The conclusion reached is \eqref{e:compactExtensionConclusion} for $\widetilde T$ with dilation $0$, 
which is \eqref{e:compactExtensionConclusion} for $T$ with dilation $i_0$.  
The only appeal to the hypothesis on $\mathcal D$ is at \eqref{e:c0}, which now reads 
\begin{equation}
\liminf_{N-M\to\infty}  \Expectation_{M\leq u<N} 
m \bigl( E \cap T ^{-b ^{i_0+i}k_u}E \cap \cdots \cap T ^{-tb ^{i_0+i}k_u}E \bigr)>0 , 
\end{equation}
with $i$ the scale \eqref{e:scale}.  This is \eqref{e:CSZdef} for $\mathcal D$, for the same Cantor set $\mathcal C$, with dilation $i_0+i$, 
and so is part of the hypothesis.  

It is worth saying why the argument is organized in this way, rather than by asserting that $(X,\mathcal D,m,\widetilde T)$ is CSZ and quoting the case $i_0=0$.  
That assertion would be false in general.  The CSZ property quantifies over all Cantor sets, and the dilations it provides for a second Cantor set 
are powers of that set's base, not of $b$.  What is true, and all that is used, is the single instance displayed above.  
\end{remark}

\subsection{Systems whose eigenvalues are roots of unity} \label{sub:rootsOfUnity}

This subsection is independent of the proof of the main theorem.  

\begin{theorem} \label{t:compactSystem} 
Let $(X,\mathcal A, m,T)$ be a measure preserving system in which the functions $h$ satisfying $T ^{q}h = h$, for some integer $q \geq 1$, 
are dense in $L^2(X)$.  Then, for every $f\in L^\infty(X)$ with $0 \leq f \leq 1$, every $\epsilon >0$, and every integer $t \geq 0$, the set of indices 
$n$ such that  
\begin{equation} \label{e:compactSystem}
\int \prod _{s=0} ^{t} T ^{s k_n } f 
\;dm \geq \int f ^{t+1}\;dm - \epsilon 
\end{equation}
is syndetic.  
\end{theorem}

The hypothesis says exactly that $T$ has pure point spectrum and that all of its eigenvalues are roots of unity.  
Indeed, for a given $q$ the functions with $T ^{q}h=h$ form a $T$-invariant subspace on which $T ^{q}$ is the identity, 
and which is therefore spanned by eigenfunctions whose eigenvalues are $q$-th roots of unity; conversely, a finite linear combination of 
eigenfunctions whose eigenvalues are roots of unity is fixed by $T ^{q}$, for $q$ the least common multiple of their orders.  
Odometers are examples.  Such a system is compact, but a compact system need not satisfy the hypothesis: 
for an irrational rotation of the circle, $T ^{q}$ is ergodic for every $q \geq 1$, so the only functions fixed by some $T ^{q}$ are the constants.  

\begin{proof}
Let $\eta \coloneqq \epsilon /(2(t+1))$, and choose $h$ and $q \geq 1$ with $T ^{q}h = h$ and $\lVert f-h\rVert _2 < \eta$.  
Let $h' \coloneqq \min (1, \max (0, \operatorname{Re} h))$.  Then $T ^{q}h' = h'$, as $h'$ is a function of $h$; 
and $\lvert f - h'\rvert \leq \lvert f-h\rvert$ pointwise, because $z \mapsto \min(1,\max(0,\operatorname{Re} z))$ is $1$-Lipschitz and fixes the values of $f$, which lie in $[0,1]$.  
Hence $\lVert f-h'\rVert _1 \leq \lVert f - h'\rVert _2 < \eta$.  

For functions $a_0 , \ldots, a_t$ and $b_0 , \ldots, b_t$ with values in $[0,1]$, telescoping gives 
$\lvert \prod_s a_s - \prod_s b_s\rvert \leq \sum_s \lvert a_s - b_s\rvert$ pointwise.  As $T$ preserves $m$, it follows that, for every $n$, 
\begin{equation}
\Bigl\lvert \int \prod _{s=0} ^{t} T ^{s k_n } f \;dm - \int \prod _{s=0} ^{t} T ^{s k_n } h' \;dm \Bigr\rvert 
\leq (t+1) \lVert f - h'\rVert _1 < \frac \epsilon 2 , 
\end{equation}
and likewise $\int h' {}^{t+1}\;dm \geq \int f ^{t+1}\;dm - \epsilon/2$.  

Let $S_q$ be the set of indices $n$ with $k_n \equiv 0 \mod q$, which is syndetic by Lemma \ref{l:syndeticResidues}.  
For $n \in S_q$ we have $T ^{sk_n}h' = h'$ for every $s$, so $\int \prod_{s=0}^t T ^{s k_n } h' \;dm = \int h' {}^{t+1}\;dm$, and therefore 
\begin{equation}
\int \prod _{s=0} ^{t} T ^{s k_n } f \;dm > \int h' {}^{t+1}\;dm - \frac \epsilon 2 \geq \int f ^{t+1}\;dm - \epsilon . 
\end{equation}
That is, \eqref{e:compactSystem} holds on the syndetic set $S_q$.  
\end{proof}
